\section{Longest Increasing Subsequence}
\label{sec:dp-lis}

\problemstatement{Given an array $\textit{arr}$ of $n$ integers, find
the length of the \textbf{longest strictly increasing subsequence}
(elements in order, but not necessarily contiguous).}

\begin{patternbox}
\emph{``Longest increasing subsequence''} is the chapter's foundational
problem: at every index, the question is whether it can \textbf{extend}
some earlier subsequence -- which depends on what the \emph{previous
chosen element's value} was, not just which index we're at.
\end{patternbox}

\subsection*{Deriving the recurrence (pick/not-pick formulation)}

Let $f(\textit{ind}, \textit{prevIndex})$ be the length of the longest
increasing subsequence achievable from index $\textit{ind}$ onward,
given that the most recently \emph{included} element (if any) is at
index $\textit{prevIndex}$ (or $-1$ if nothing has been included yet).
At each index, either skip it ($f(\textit{ind}{+}1, \textit{prevIndex})$),
or include it -- only valid if $\textit{prevIndex}=-1$ or
$\textit{arr}[\textit{prevIndex}] < \textit{arr}[\textit{ind}]$ -- gaining
$1 + f(\textit{ind}{+}1, \textit{ind})$:
\[
  f(\textit{ind}, \textit{prevIndex}) = \max\Big(
  f(\textit{ind}{+}1, \textit{prevIndex}),\ \textit{canTake} \Big), \quad
  \textit{canTake} = 1 + f(\textit{ind}{+}1, \textit{ind})\ \text{(if valid)}.
\]

\subsection*{Approach 1: Recursion}

\begin{lstlisting}
#include <bits/stdc++.h>
using namespace std;

int lisRecursive(int ind, int prevIndex, vector<int>& arr) {
    int n = arr.size();
    if (ind == n) return 0;

    int notTake = lisRecursive(ind + 1, prevIndex, arr);
    int take = 0;
    if (prevIndex == -1 || arr[prevIndex] < arr[ind]) {
        take = 1 + lisRecursive(ind + 1, ind, arr);
    }
    return max(notTake, take);
}

int main() {
    vector<int> arr = {10, 9, 2, 5, 3, 7, 101, 18};
    cout << lisRecursive(0, -1, arr) << endl; // 4
    return 0;
}
\end{lstlisting}

\begin{complexitybox}[Approach 1 Complexity]
\textbf{Time.} $O(2^n)$.
\textbf{Space.} $O(n)$ recursion stack.
\end{complexitybox}

\subsection*{Approach 2: Memoization}

\begin{lstlisting}
#include <bits/stdc++.h>
using namespace std;

int lisMemo(int ind, int prevIndex, vector<int>& arr, vector<vector<int>>& memo) {
    int n = arr.size();
    if (ind == n) return 0;
    int prevShift = prevIndex + 1; // shift by 1 so -1 maps to 0
    if (memo[ind][prevShift] != -1) return memo[ind][prevShift];

    int notTake = lisMemo(ind + 1, prevIndex, arr, memo);
    int take = 0;
    if (prevIndex == -1 || arr[prevIndex] < arr[ind]) {
        take = 1 + lisMemo(ind + 1, ind, arr, memo);
    }
    return memo[ind][prevShift] = max(notTake, take);
}

int main() {
    vector<int> arr = {10, 9, 2, 5, 3, 7, 101, 18};
    int n = arr.size();
    vector<vector<int>> memo(n, vector<int>(n + 1, -1));
    cout << lisMemo(0, -1, arr, memo) << endl; // 4
    return 0;
}
\end{lstlisting}

\begin{complexitybox}[Approach 2 Complexity]
\textbf{Time.} $O(n^2)$.
\textbf{Space.} $O(n^2)$ memo $+$ $O(n)$ recursion stack.
\end{complexitybox}

\subsection*{Approach 3: Tabulation (Ending-Here Reformulation)}

\begin{ideabox}[Why Tabulation Uses a Cleaner, Different-Looking Table]
Directly tabulating the $(\textit{ind}, \textit{prevIndex})$ recursion
is awkward (the shifted-index bookkeeping is unpleasant to iterate over
bottom-up). The standard simplification: define
$\textit{dp}[i]$ = length of the LIS \emph{ending exactly at} index $i$.
This is mathematically equivalent information, just reorganized; it
gives a much cleaner loop and is how this problem is conventionally
tabulated.
\end{ideabox}

\[
  \textit{dp}[i] = 1 + \max_{j<i,\ \textit{arr}[j]<\textit{arr}[i]}
  \textit{dp}[j] \quad (\text{or } 1 \text{ if no such } j \text{ exists}).
\]

\begin{lstlisting}
#include <bits/stdc++.h>
using namespace std;

int lisTabulation(vector<int>& arr) {
    int n = arr.size();
    vector<int> dp(n, 1); // every element alone is an LIS of length 1

    int best = 1;
    for (int i = 0; i < n; i++) {
        for (int j = 0; j < i; j++) {
            if (arr[j] < arr[i]) {
                dp[i] = max(dp[i], 1 + dp[j]);
            }
        }
        best = max(best, dp[i]);
    }
    return best;
}

int main() {
    vector<int> arr = {10, 9, 2, 5, 3, 7, 101, 18};
    cout << lisTabulation(arr) << endl; // 4
    return 0;
}
\end{lstlisting}

\subsection*{Dry run of the tabulation table}

\dryrun{Take $\textit{arr} = [10,9,2,5,3,7,101,18]$.}

\begin{center}
\begin{tabular}{c|c|c|c|c|c|c|c|c}
$i$ & 0 & 1 & 2 & 3 & 4 & 5 & 6 & 7 \\ \hline
$\textit{arr}[i]$ & 10 & 9 & 2 & 5 & 3 & 7 & 101 & 18 \\ \hline
$\textit{dp}[i]$ & 1 & 1 & 1 & 2 & 2 & 3 & 4 & 4
\end{tabular}
\end{center}

$\textit{dp}[5]=3$ comes from extending the chain ending at index $3$
(value $5$, $\textit{dp}=2$): $2,5,7$. $\textit{dp}[6]=4$ extends that
same chain with $101$: $2,5,7,101$. The maximum over the whole table is
$4$.

\begin{complexitybox}[Approach 3 Complexity]
\textbf{Time.} $O(n^2)$ -- a nested loop comparing every pair.
\textbf{Space.} $O(n)$ for the \textit{dp} array -- already minimal;
this problem's natural tabulated form has no further space reduction
the way 2D-table problems in earlier chapters did.
\end{complexitybox}

\begin{ideabox}[A Faster $O(n \log n)$ Algorithm Exists]
Section~\ref{sec:dp-lis-binary-search} develops a completely different
$O(n \log n)$ technique for this same problem, using binary search
instead of dynamic programming's pairwise comparison -- worth knowing
as the production-quality solution once the $O(n^2)$ DP version here is
understood.
\end{ideabox}

\begin{takeawaybox}
``Ending-here'' tables answer their question with a \textbf{maximum
over the entire filled table}, not a single corner cell -- the first
new structural lesson of this chapter, and the one to watch for in
every section that follows.
\end{takeawaybox}
